
\documentclass[11pt,letterpaper]{article}

% ==========================================
% PAQUETES
% ==========================================
\usepackage[utf8]{inputenc}
\usepackage[T1]{fontenc}
\usepackage[spanish,es-tabla]{babel}
\usepackage{lmodern}
\usepackage{amsmath,amssymb}
\usepackage{geometry}
\usepackage{xcolor}
\usepackage{tikz}
\usepackage[most]{tcolorbox}
\usepackage{fancyhdr}
\usepackage{enumitem}
\usepackage{array}
\usepackage{lastpage}

% ==========================================
% TIKZ: ANTES DE BEGIN DOCUMENT
% ==========================================
\usetikzlibrary{
 arrows.meta,
 calc,
 positioning,
 patterns,
 babel
}

% ==========================================
% CONFIGURACIÓN DE PÁGINA
% ==========================================
\geometry{
 left=1.8cm,
 right=1.8cm,
 top=2.3cm,
 bottom=2cm
}

\setlength{\parindent}{0pt}
\setlength{\parskip}{5pt}
\setlength{\headheight}{15pt}

% ==========================================
% COLORES
% ==========================================
\definecolor{azul}{RGB}{20,65,115}
\definecolor{celeste}{RGB}{232,243,252}
\definecolor{verde}{RGB}{0,130,100}
\definecolor{morado}{RGB}{110,50,160}
\definecolor{gris}{RGB}{246,248,251}
\definecolor{naranja}{RGB}{230,110,35}

% ==========================================
% ENCABEZADO Y PIE DE PÁGINA
% ==========================================
\pagestyle{fancy}
\fancyhf{}

\lhead{
 \small\textcolor{azul}{
 \textbf{Física -- Educación Diversificada}
 }
}

\rhead{
 \small\textcolor{azul}{
 Práctica de gráficas
 }
}

\cfoot{
 \small Página \thepage\ de \pageref{LastPage}
}

\renewcommand{\headrulewidth}{0.5pt}

% ==========================================
% RECUADROS
% ==========================================
\newtcolorbox{seccion}[1]{
 enhanced,
 breakable,
 colback=celeste,
 colframe=azul,
 colbacktitle=azul,
 coltitle=white,
 title={#1},
 fonttitle=\bfseries,
 arc=2mm,
 boxrule=0.8pt
}

\newtcolorbox{ejercicio}[1]{
 enhanced,
 breakable,
 colback=white,
 colframe=azul!55,
 colbacktitle=celeste,
 coltitle=azul,
 title={Ejercicio #1},
 fonttitle=\bfseries,
 arc=2mm,
 boxrule=0.6pt
}

\newtcolorbox{reto}{
 enhanced,
 breakable,
 colback=orange!5,
 colframe=naranja,
 colbacktitle=naranja,
 coltitle=white,
 title={Desafío final},
 fonttitle=\bfseries,
 arc=2mm
}

% ==========================================
% COMANDOS
% ==========================================
\newcommand{\respuesta}{
 \par\vspace{0.3cm}
 \textcolor{morado}{\textbf{Respuesta:}}
 \hrulefill
 \par\vspace{0.3cm}
}

\newcommand{\espacio}[1]{
 \par\vspace{#1}
}

\newcommand{\casilla}{
 $\square$
}

% ==========================================
% GRÁFICA POSICIÓN-TIEMPO
% ==========================================
% #1 = lista de coordenadas
%
% Ejemplo:
% \grafpos{(0,0) (10,20) (20,30)}
%
% El eje vertical representa posición d.
% ==========================================
\newcommand{\grafpos}[1]{
\begin{tikzpicture}[
 x=0.12cm,
 y=0.085cm,
 >=Stealth
]

% Cuadrícula
\draw[
 step=10,
 gray!20,
 very thin
]
(0,0) grid (50,40);

% Eje horizontal
\draw[->,thick]
(0,0)--(55,0)
node[right] {$t$ (s)};

% Eje vertical
\draw[->,thick]
(0,0)--(0,46)
node[above] {$d$ (m)};

% Escala horizontal
\foreach \x in {10,20,30,40,50}{
 \draw (\x,0)--(\x,-1);
 \node[
 below,
 font=\scriptsize
 ] at (\x,-1) {\x};
}

% Escala vertical
\foreach \y in {10,20,30,40}{
 \draw (0,\y)--(-1,\y);
 \node[
 left,
 font=\scriptsize
 ] at (-1,\y) {\y};
}

% Trayectoria
\draw[
 very thick,
 azul
]
plot coordinates {#1};

\end{tikzpicture}
}

% ==========================================
% GRÁFICA RAPIDEZ-TIEMPO
% ==========================================
\newcommand{\grafvel}[1]{
\begin{tikzpicture}[
 x=0.65cm,
 y=0.30cm,
 >=Stealth
]

\draw[
 step=1,
 gray!15,
 very thin
]
(0,0) grid (10,12);

\draw[->,thick]
(0,0)--(11,0)
node[right] {$t$ (s)};

\draw[->,thick]
(0,0)--(0,13)
node[above] {$v$ (m/s)};

\foreach \x in {2,4,6,8,10}{
 \node[
 below,
 font=\scriptsize
 ] at (\x,0) {\x};
}

\foreach \y in {2,4,6,8,10,12}{
 \node[
 left,
 font=\scriptsize
 ] at (0,\y) {\y};
}

\draw[
 very thick,
 azul
]
plot coordinates {#1};

\end{tikzpicture}
}

% ==========================================
% DOCUMENTO
% ==========================================
\begin{document}

% ==========================================
% PORTADA DE LA PRÁCTICA
% ==========================================
\begin{center}

{\Large
\bfseries
\textcolor{azul}{
PRÁCTICA DE FÍSICA
}
}

\vspace{0.2cm}

{\large
Gráficas del movimiento rectilíneo
}

\vspace{0.2cm}

\textcolor{verde}{
\textbf{
Pendientes, velocidades,
distancias y desplazamientos
}
}

\end{center}

\begin{tcolorbox}[
 colback=gris,
 colframe=azul,
 arc=2mm
]

\textbf{Estudiante:}
\hrulefill

\vspace{0.3cm}

\textbf{Sección:}
\underline{\hspace{3cm}}
\hfill
\textbf{Fecha:}
\underline{\hspace{3cm}}

\end{tcolorbox}

\begin{tcolorbox}[
 colback=celeste,
 colframe=azul,
 title={Indicaciones}
]

\begin{itemize}[leftmargin=*]
\item Analice cuidadosamente cada gráfica.
\item Identifique los intervalos de tiempo.
\item Muestre los procedimientos matemáticos.
\item Utilice las unidades correspondientes.
\item Redondee a dos decimales si es necesario.
\item Interprete los resultados obtenidos.
\end{itemize}

\end{tcolorbox}

% ==========================================
% SECCIÓN I
% ==========================================
\section*{I. Interpretación de gráficas}

\begin{seccion}{
Gráficas posición--tiempo
}
Observe cada gráfica y responda
las preguntas planteadas.

En estas gráficas, $d$ representa
la posición respecto al origen.
\end{seccion}
\newpage
% EJERCICIO 1
\begin{ejercicio}{1}

La siguiente gráfica representa
la posición de un estudiante
durante su recorrido.

\begin{center}
\grafpos{
(0,0)
(10,20)
(20,20)
(30,40)
(40,10)
}
\end{center}

Determine:

\begin{enumerate}[label=\alph*)]
\item Posición inicial.
\item Posición a los $10\,\mathrm{s}$.
\item Posición a los $20\,\mathrm{s}$.
\item Intervalo en que permanece en reposo.
\item Intervalo en que regresa hacia el origen.
\end{enumerate}

\espacio{2cm}
\respuesta

\end{ejercicio}

\newpage

% EJERCICIO 2
\begin{ejercicio}{2}

La gráfica representa
el movimiento de un ciclista.

\begin{center}
\grafpos{
(0,10)
(10,30)
(20,30)
(30,0)
(40,20)
}
\end{center}

Determine:

\begin{enumerate}[label=\alph*)]
\item Posición inicial.
\item Posición final.
\item Intervalos de movimiento positivo.
\item Intervalos de movimiento negativo.
\item Intervalos de reposo.
\end{enumerate}

\espacio{2cm}
\respuesta

\end{ejercicio}
\newpage
% ==========================================
% SECCIÓN II
% ==========================================
\section*{II. Pendiente y velocidad}

\begin{seccion}{
Cálculo de velocidades
}

Determine la velocidad de cada tramo
utilizando la pendiente.

\[
v=\frac{d_f-d_i}{t_f-t_i}
\]

\end{seccion}

% EJERCICIO 3
\begin{ejercicio}{3}

Observe la gráfica:

\begin{center}
\grafpos{
(0,0)
(10,30)
(20,30)
(30,10)
(40,0)
}
\end{center}

Calcule:

\begin{enumerate}[label=\alph*)]
\item Velocidad entre $0$ y $10\,\mathrm{s}$.
\item Velocidad entre $10$ y $20\,\mathrm{s}$.
\item Velocidad entre $20$ y $30\,\mathrm{s}$.
\item Velocidad entre $30$ y $40\,\mathrm{s}$.
\end{enumerate}

\espacio{3cm}
\respuesta

\end{ejercicio}

\newpage

% EJERCICIO 4
\begin{ejercicio}{4}

Un automóvil presenta
el siguiente movimiento.

\begin{center}
\grafpos{
(0,0)
(10,10)
(20,40)
(30,40)
(40,20)
}
\end{center}

Determine:

\begin{enumerate}[label=\alph*)]
\item Velocidad de cada intervalo.
\item Intervalo de mayor rapidez.
\item Intervalo de reposo.
\item Intervalo de velocidad negativa.
\end{enumerate}

\espacio{3cm}
\respuesta

\end{ejercicio}
\newpage
% ==========================================
% SECCIÓN III
% ==========================================
\section*{III. Distancia y desplazamiento}

\begin{seccion}{
Distancia total y desplazamiento
}

Calcule la distancia total recorrida
y el desplazamiento del móvil.

Recuerde:

\[
\Delta d=d_f-d_i
\]

La distancia total es la suma
de las longitudes recorridas
en cada intervalo.

\end{seccion}

% EJERCICIO 5
\begin{ejercicio}{5}

Observe la gráfica:

\begin{center}
\grafpos{
(0,0)
(10,30)
(20,10)
(30,40)
(40,20)
}
\end{center}

Determine:

\begin{enumerate}[label=\alph*)]
\item Distancia recorrida en cada intervalo.
\item Distancia total.
\item Desplazamiento.
\item Rapidez media.
\item Velocidad media.
\end{enumerate}

\espacio{3cm}
\respuesta

\end{ejercicio}

\newpage

% EJERCICIO 6
\begin{ejercicio}{6}

Un estudiante se desplaza
según la siguiente gráfica.

\begin{center}
\grafpos{
(0,20)
(10,40)
(20,10)
(30,10)
(40,30)
}
\end{center}

Determine:

\begin{enumerate}[label=\alph*)]
\item Posición inicial y final.
\item Distancia total recorrida.
\item Desplazamiento.
\item Rapidez media.
\item Velocidad media.
\end{enumerate}

\espacio{3cm}
\respuesta

\end{ejercicio}

% ==========================================
% SECCIÓN IV
% ==========================================
\newpage

\section*{IV. Gráficas rapidez--tiempo}

\begin{seccion}{
Cálculo de distancias mediante áreas
}

Calcule las áreas geométricas
correspondientes a cada intervalo.

\[
D=\text{Área bajo la gráfica}
\]

Utilice las fórmulas de área
del triángulo, rectángulo
y trapecio.

\end{seccion}

% EJERCICIO 7
\begin{ejercicio}{7}

La rapidez de un automóvil
varía según la gráfica.

\begin{center}
\grafvel{
(0,0)
(2,6)
(5,6)
(8,0)
}
\end{center}

Determine:

\begin{enumerate}[label=\alph*)]
\item Distancia entre $0$ y $2\,\mathrm{s}$.
\item Distancia entre $2$ y $5\,\mathrm{s}$.
\item Distancia entre $5$ y $8\,\mathrm{s}$.
\item Distancia total recorrida.
\end{enumerate}

\espacio{3cm}
\respuesta

\end{ejercicio}

\newpage

% EJERCICIO 8
\begin{ejercicio}{8}

Un motociclista presenta
la siguiente gráfica rapidez--tiempo.

\begin{center}
\grafvel{
(0,0)
(2,4)
(4,4)
(6,10)
(8,10)
(10,0)
}
\end{center}

Determine:

\begin{enumerate}[label=\alph*)]
\item Distancia de $0$ a $2\,\mathrm{s}$.
\item Distancia de $2$ a $4\,\mathrm{s}$.
\item Distancia de $4$ a $6\,\mathrm{s}$.
\item Distancia de $6$ a $8\,\mathrm{s}$.
\item Distancia de $8$ a $10\,\mathrm{s}$.
\item Distancia total.
\end{enumerate}

\espacio{3cm}
\respuesta

\end{ejercicio}
\newpage
% EJERCICIO 9
\begin{ejercicio}{9}

La rapidez de un vehículo
se representa mediante los puntos:

\[
(0,0),\quad
(3,9),\quad
(6,9),\quad
(9,0)
\]

donde el tiempo está en segundos
y la rapidez en $\mathrm{m/s}$.

Los puntos se unen mediante
segmentos rectos.

\begin{enumerate}[label=\alph*)]
\item Construya la gráfica.
\item Identifique las figuras geométricas.
\item Calcule la distancia total.
\item Determine la rapidez media.
\end{enumerate}

\espacio{5cm}
\respuesta

\end{ejercicio}

% ==========================================
% SECCIÓN V
% ==========================================
\newpage

\section*{V. Gráficas velocidad--tiempo}

\begin{seccion}{
Aceleración y desplazamiento
}

Utilice la pendiente para calcular
la aceleración:

\[
a=\frac{v_f-v_i}{t_f-t_i}
\]

El área con signo representa
el desplazamiento.

La suma de los valores absolutos
de las áreas representa la
distancia total recorrida.

\end{seccion}

% EJERCICIO 10
\begin{ejercicio}{10}

La velocidad de un automóvil
aumenta uniformemente
de $0$ a $12\,\mathrm{m/s}$
durante $6\,\mathrm{s}$.

\begin{enumerate}[label=\alph*)]
\item Construya la gráfica velocidad--tiempo.
\item Determine la aceleración.
\item Calcule el desplazamiento.
\end{enumerate}

\espacio{5cm}
\respuesta

\end{ejercicio}
\newpage
% EJERCICIO 11
\begin{ejercicio}{11}

Un automóvil disminuye
su velocidad uniformemente
de $20\,\mathrm{m/s}$
a $0\,\mathrm{m/s}$
en $5\,\mathrm{s}$.

Determine:

\begin{enumerate}[label=\alph*)]
\item La aceleración.
\item El desplazamiento durante el frenado.
\item La gráfica velocidad--tiempo.
\end{enumerate}

\espacio{5cm}
\respuesta

\end{ejercicio}

\newpage

% EJERCICIO 12
\begin{ejercicio}{12}

La velocidad de un móvil
cambia de la siguiente manera:

\begin{center}
\begin{tikzpicture}[
 x=0.8cm,
 y=0.35cm,
 >=Stealth
]

% Cuadrícula
\draw[
 step=1,
 gray!15,
 very thin
]
(0,-5) grid (8,7);

% Ejes
\draw[->,thick]
(0,0)--(9,0)
node[right] {$t$ (s)};

\draw[->,thick]
(0,-5)--(0,8)
node[above] {$v$ (m/s)};

% Gráfica
\draw[very thick,azul]
(0,6)--(2,6)--(4,0)--(6,-4)--(8,-4);

% Escala horizontal
\foreach \x in {2,4,6,8}{
 \draw (\x,0)--(\x,-0.15);
 \node[
 below right,
 font=\scriptsize
 ] at (\x,0) {\x};
}

% Escala vertical
\foreach \y in {-4,-2,2,4,6}{
 \draw (0,\y)--(-0.15,\y);
 \node[
 left,
 font=\scriptsize
 ] at (0,\y) {\y};
}

\end{tikzpicture}
\end{center}

Determine:

\begin{enumerate}[label=\alph*)]
\item Aceleración de cada intervalo.
\item Desplazamiento positivo.
\item Desplazamiento negativo.
\item Desplazamiento total.
\item Distancia total recorrida.
\end{enumerate}

\espacio{3cm}
\respuesta

\end{ejercicio}

% ==========================================
% SECCIÓN VI
% ==========================================
\newpage

\section*{VI. Problemas integradores}

\begin{seccion}{
Aplicación de los conceptos
}

Resuelva los siguientes problemas
mostrando todos los procedimientos.

\end{seccion}

% EJERCICIO 13
\begin{ejercicio}{13}

Un ciclista inicia su recorrido
en la posición $d=0\,\mathrm{m}$.

Durante los primeros $10\,\mathrm{s}$
avanza hasta $40\,\mathrm{m}$.

Permanece detenido durante
$5\,\mathrm{s}$.

Finalmente regresa hasta
la posición $10\,\mathrm{m}$
en otros $10\,\mathrm{s}$.

Determine:

\begin{enumerate}[label=\alph*)]
\item Gráfica posición--tiempo.
\item Velocidad de cada intervalo.
\item Distancia total.
\item Desplazamiento.
\item Rapidez media.
\item Velocidad media.
\end{enumerate}

\espacio{5cm}
\respuesta

\end{ejercicio}
\newpage
% EJERCICIO 14
\begin{ejercicio}{14}

Un vehículo parte del reposo
y aumenta uniformemente su rapidez
hasta $16\,\mathrm{m/s}$
durante $8\,\mathrm{s}$.

Mantiene esa rapidez durante
$5\,\mathrm{s}$ y después
frena uniformemente hasta detenerse
en $4\,\mathrm{s}$.

Determine:

\begin{enumerate}[label=\alph*)]
\item Gráfica rapidez--tiempo.
\item Distancia durante la aceleración.
\item Distancia a rapidez constante.
\item Distancia durante el frenado.
\item Distancia total recorrida.
\end{enumerate}

\espacio{5cm}
\respuesta

\end{ejercicio}

\newpage

% EJERCICIO 15
\begin{ejercicio}{15}

Un automóvil se desplaza
según la siguiente gráfica.

\begin{center}
\grafpos{
(0,0)
(10,40)
(20,40)
(30,10)
(40,30)
}
\end{center}

Determine:

\begin{enumerate}[label=\alph*)]
\item Velocidad de cada tramo.
\item Distancia total.
\item Desplazamiento.
\item Rapidez media.
\item Velocidad media.
\item Intervalo de mayor rapidez.
\end{enumerate}

\espacio{4cm}
\respuesta

\end{ejercicio}

% ==========================================
% SECCIÓN VII
% ==========================================
\newpage

\section*{VII. Desafío final}

\begin{reto}

\textbf{Ejercicio 16.}

Un autobús inicia su movimiento
desde el reposo.

Durante los primeros
$5\,\mathrm{s}$ aumenta
uniformemente su rapidez
hasta $10\,\mathrm{m/s}$.

Luego mantiene esa rapidez
durante $8\,\mathrm{s}$.

Finalmente disminuye
uniformemente su rapidez
hasta detenerse en $5\,\mathrm{s}$.

Determine:

\begin{enumerate}[label=\alph*)]
\item Gráfica rapidez--tiempo.
\item Aceleración inicial.
\item Aceleración durante el frenado.
\item Distancia de cada intervalo.
\item Distancia total.
\item Rapidez media de todo el recorrido.
\end{enumerate}

\end{reto}

\espacio{7cm}
\respuesta

% ==========================================
% AUTOEVALUACIÓN
% ==========================================
\newpage

\section*{Autoevaluación}

\begin{tcolorbox}[
 colback=celeste,
 colframe=azul,
 title={Revisión de mi aprendizaje}
]

Marque las habilidades que domina.

\vspace{0.3cm}

\casilla\ Interpreto gráficas
posición--tiempo.

\vspace{0.4cm}

\casilla\ Calculo velocidades
mediante pendientes.

\vspace{0.4cm}

\casilla\ Identifico intervalos
de reposo y retroceso.

\vspace{0.4cm}

\casilla\ Calculo distancia total
y desplazamiento.

\vspace{0.4cm}

\casilla\ Calculo distancias
mediante áreas geométricas.

\vspace{0.4cm}

\casilla\ Interpreto gráficas
velocidad--tiempo.

\vspace{0.4cm}

\casilla\ Determino aceleraciones
mediante pendientes.

\end{tcolorbox}

\vspace{0.5cm}

\textbf{¿Qué tema necesito practicar más?}

\espacio{2cm}
\hrulefill

\vspace{0.5cm}

\textbf{Observaciones del docente:}

\espacio{2cm}
\hrulefill

\end{document}
